%% abstract.tex: abstract in PT and EN  (FEUP regulations)
%% -------------------------------------------------------
%

\chapter*{Resumo}
\hspace{\parindent}
Esta dissertação aborda o planeamento de caminhos multi-robô, \emph{Multi-Robot Path Planning} (MRPP), para frotas de robôs móveis autónomos
que operam em ambientes partilhados. Neste contexto, vários robôs devem alcançar
os seus objetivos, evitando colisões e respeitando restrições temporais de
coordenação. O trabalho foca-se na aplicação de MRPP a uma arquitetura
centralizada de coordenação de frotas multi-robô, onde as decisões de
planeamento devem ser computacionalmente eficientes e compatíveis com restrições
de execução prática.

O trabalho tem por base a arquitetura de coordenação desenvolvida no Instituto
de Engenharia de Sistemas e Computadores, Tecnologia e Ciência (INESC TEC), no
âmbito do Centro de Robótica na Indústria e Sistemas Inteligentes (CRIIS) e do
\emph{Industry and Innovation Laboratory} (iiLab). Esta arquitetura utiliza um
planeador sequencial \emph{Time-Enhanced A*} (TEA*), no qual os robôs são
planeados um de cada vez segundo uma ordenação de prioridades. Embora esta
abordagem seja computacionalmente eficiente e adequada a cenários de aplicação
prática, o seu desempenho é altamente sensível à ordenação escolhida, sendo
inviável avaliar exaustivamente todas as possíveis permutações para frotas de
maior dimensão.

Para mitigar esta limitação, esta dissertação desenvolve um conjunto de
heurísticas de prioritização que geram ordenações alternativas dos robôs com
base em distância, topologia, distribuição espacial, conflitos previstos e
informação de replaneamento. Estas ordenações são primeiro avaliadas num
ambiente de teste isolado do planeador e, posteriormente, integradas numa
abordagem de planeamento paralelo, onde várias instâncias do TEA* são executadas
em paralelo sob um tempo máximo de planeamento. O objetivo não é substituir o
planeador existente, mas avaliar se a sua robustez e a qualidade das soluções
podem ser melhoradas através da exploração de várias ordenações promissoras no
mesmo ciclo de planeamento.

Os resultados experimentais mostram que as sequências de robôs geradas por
heurísticas melhoraram a fiabilidade na obtenção de soluções válidas pelo TEA*
nos cenários restritos e de grande escala avaliados, quando comparadas com a
seleção aleatória de prioridades. A abordagem de planeamento paralelo aumentou
ainda, nos cenários avaliados, a capacidade de identificar ordenações úteis
dentro do tempo de planeamento disponível, reduzindo a dependência de uma única
permutação de prioridades.

Relativamente à execução integrada, a análise da arquitetura completa de
coordenação do CRIIS identifica operações redundantes de replaneamento e
processamento temporalmente inconsistente de pedidos de planeamento e
atualizações de estado dos robôs, motivando a introdução de mecanismos de
agrupamento destinados a melhorar a consistência da execução. No geral, esta
dissertação estende a arquitetura centralizada de coordenação de frotas através
de heurísticas de prioritização, planeamento TEA* paralelo e mecanismos de
agrupamento, com melhorias observadas no contexto de MRPP avaliado.

\vfill
\noindent\textbf{Palavras-chave:} Planeamento de Caminhos Multi-Robô, Coordenação Centralizada de
Frotas, Heurísticas de Priorização, Time-Enhanced A*, Planeamento Paralelo, Coordenação de Robôs.


\chapter*{Abstract}
\hspace{\parindent}
This dissertation addresses centralized {Multi-Robot Path Planning} (MRPP) for fleets of autonomous mobile
robots operating in shared environments. In this context, multiple robots must
reach their assigned goals while avoiding collisions and satisfying temporal
coordination constraints. The work focuses on the application of MRPP to a
centralized multi-robot fleet coordination framework, where planning decisions
must remain computationally efficient and compatible with practical execution
constraints.

The work builds on the existing coordination stack developed at the Institute
for Systems and Computer Engineering, Technology and Science (INESC TEC),
within the Centre for Robotics in Industry and Intelligent Systems (CRIIS) and
the Industry and Innovation Laboratory (iiLab). This framework relies on a
sequential {Time-Enhanced A*} (TEA*) planner, where robots are planned one
at a time according to a priority ordering. While this approach is
computationally efficient and well suited for practical deployment, its
performance is highly sensitive to the selected ordering, and exhaustive
evaluation of all possible permutations is infeasible for large fleets.

To address this limitation, this dissertation develops a set of prioritization
heuristics that generate alternative robot orderings using distance, topology,
surroundings, predicted conflicts, and replanning feedback. These orderings are
first evaluated in an isolated planner testing framework and are later
integrated into a parallel planning approach, where multiple TEA* instances are
executed in parallel under a fixed planning time budget. The objective is not to
replace the existing planner, but to evaluate whether its robustness and solution quality can be
improved by exploring several promising robot orderings within the same planning
cycle.

The experimental results show that heuristic-generated robot sequences improved the
valid-solution reliability of TEA* in the evaluated constrained and large-scale scenarios when
compared with random priority selection, especially in scenarios where the ordering has a strong impact on solution
feasibility. The parallel planner further increased the ability to identify effective ordering
alternatives within the available planning time, reducing the dependence on a single priority
permutation.

Regarding integrated execution, the analysis of the complete CRIIS coordination
stack identifies redundant replanning operations and temporally inconsistent
processing of planning requests and robot state updates, which motivated batching
mechanisms aimed at improving execution consistency during integrated operation. Overall, the dissertation extends
the existing centralized fleet coordination framework through prioritization
heuristics, parallel TEA* planning, and batching mechanisms, with observed improvements in the
evaluated MRPP setting.

\vfill

\noindent\textbf{Keywords:} Multi-Robot Path Planning, Centralized Fleet Coordination,
Prioritization Heuristics, Time-Enhanced A*, Parallel Planning, Robot Coordination.